\section{Related Work}

\paragraph{POMDP planning and belief-dependent rewards.}
POMDPs formalize decision-making under state uncertainty \citep{smallwood1973,kaelbling1998}. Point-based offline solvers approximate the value function on reachable beliefs, including PBVI, HSVI, and SARSOP \citep{pineau2003,smith2004,kurniawati2008,shani2013}. Online solvers plan from the current belief. POMCP uses Monte Carlo tree search with UCB1 action selection and rollout evaluation \citep{silver2010}, DESPOT searches regularized sparse belief trees \citep{ye2017}, and POMCPOW uses progressive widening for continuous spaces \citep{sunberg2018}. Value of Information Monte Carlo planning selectively skips observation branches with low expected decision value \citep{laouar2026voi}. Our planner instead adds expected information gain to the objective. These are different uses of information within planning, and our comparisons do not establish superiority over state-of-the-art POMDP solvers generally.

\citet{araya2010} introduced $\rho$-POMDPs with belief-dependent rewards. Their $\rho$ denotes the whole reward, whereas we reserve it for an information term added to the ordinary task reward. Convex belief rewards support convex value functions and piecewise-linear convex approximations. Expected information gain is concave in the belief, so that machinery does not apply directly. \citet{fehr2018} treat non-convex rewards under a Lipschitz hypothesis, and \citet{benchetrit2025} develop an anytime planner for continuous spaces. For a finite observation model, strictly positive likelihoods make expected information gain Lipschitz on the belief simplex. Zero likelihoods, including noiseless sensing, can violate this condition. Our exact tree recursion does not rely on that approximation guarantee.

POMDPs with information rewards were developed for active cooperative perception \citep{spaan2015}, and their relation to $\rho$-POMDPs is discussed by \citet{satsangi2018}. Earlier preference-elicitation work selected queries by sequential value of information while retaining state-dependent rewards \citep{boutilier2002}. These approaches establish the motivation for valuing information explicitly. Our focus is the coefficient on that value and its relation to expected sensing usage.

\paragraph{Information value and experimental design.}
Classical value of information measures the improvement an observation can bring to a decision \citep{howard1966}. Bayesian experimental design instead commonly selects observations by expected information gain \citep{lindley1956,foster2021dad}. The criteria coincide under a particular reward model: when the reward is a proper log score, the expected improvement in score is the expected information gain \citep{bernardo1979}. This supplies the exact reward--information identification used here. Sequential experimental design is also the closest structural analogue to our observe-then-commit setting.

Information-directed sampling (IDS) trades instantaneous regret against information gain through their ratio, rather than through a fixed additive coefficient \citep{russo2014}. We include an observe-then-commit IDS baseline in Appendix~\ref{app:ids}. \citet{walraven2024} connect active inference with information gathering in POMDPs. Our contribution within this setting is to define and evaluate a calibration procedure for the weight in a specified $\rho$-POMDP policy family.

\paragraph{Active inference and the EFE formulation.}
Active inference casts perception and action as variational inference \citep{friston2010,parr2022}. EFE combines pragmatic and epistemic terms \citep{friston2015,parr2019,dacosta2020}. Its recursive sophisticated-inference form conditions future choices on future beliefs \citep{friston2021}. \citet{dacosta2020b} prove reward-Bellman optimality for finite-horizon recursive inference in the limit of precise reward-encoding preferences, with exact state estimation. That limiting result does not imply that a fixed finite information weight is reward-optimal.

Several distinct constructions occur under the EFE name. \citet{millidge2021} distinguish EFE from a naive future extension of variational free energy and propose Free Energy of the Expected Future (FEEF). \citet{champion2024} examine which EFE expressions follow from common root definitions and the restrictions those roots impose on prior preferences. \citet{devries2025} derive EFE-based planning from a model augmented with preference and epistemic priors. The generalised free energy of \citet{parr2019} is another distinct functional. Our equivalence concerns only the recursive, reward-based objective defined in Section~\ref{sec:methodology}, with expected information gain about the hidden state as its epistemic term. It does not assert an unchanged coefficient for these alternative constructions. Broader comparisons with decision-theoretic objectives include an EFE optimality-gap analysis \citep{wei2024voi} and a survey of formal equivalences \citep{sweeney2026equivalences}.

Scaling work combines active inference with deep models, reinforcement-learning objectives, and Monte Carlo tree search \citep{fountas2020,tschantz2020,maisto2025}. Our MCTS-EFE variant follows this direction. Its controls test exploration constants, rollout choices, and two internal components, while leaving exact beliefs, exact commit values, and EFE leaf evaluation coupled (Appendix~\ref{app:pomcp}).

\paragraph{Exploration and control as inference.}
Control-as-inference relates maximum-entropy control to probabilistic inference \citep{todorov2007,levine2018}, with soft actor-critic and KL-minimizing control as algorithmic instances \citep{haarnoja2018,rawlik2012}. Its treatment of value differs from active inference \citep{millidge2020}. EFE can produce intrinsic-motivation behavior \citep{sajid2021}, while curiosity, count-based exploration, VIME, and random network distillation use weighted exploration bonuses \citep{schmidhuber1991,pathak2017,bellemare2016,houthooft2016,burda2019}. \citet{oudeyer2007} survey these intrinsic signals. Bayes-adaptive planning and Bayesian reinforcement learning primarily address uncertainty about the model \citep{duff2002,guez2013,ghavamzadeh2015}. We study uncertainty about state under a known model, and retain dependence on a shared reward convention when describing $w=1$ as untuned.

\paragraph{Budgets and dual control.}
Rational inattention constrains or prices information capacity \citep{sims2003,matejka2015}. Constrained MDPs use Lagrange multipliers on expected resource costs \citep{altman1999}, and constrained POMDP methods track admissible cost alongside belief \citep{kim2011cpomdp}. Automatic temperature tuning in soft actor-critic drives policy entropy toward a target \citep{haarnoja2018sac}. These provide close conceptual precedents for resource control. Our weight multiplies information gain, however, while the target concerns observation count or cost. It is not generally the multiplier of a sensing-cap constraint. Section~\ref{sec:budget_positioning} gives the detailed distinctions, including the difference from a dual formulation of EFE itself \citep{kouw2026bethe}. The contribution is an operational calibration for this policy family, its exact scale transformation under stated assumptions, and explicit treatment of finite-grid crossings and usage gaps.
